% answers_short.tex —— 2005 年·普通高等学校招生考试·福建卷·理科 · 简短答案速查
% 由 tools/make_short_answers.py 自 answers.tex 的 \gksolhead/\gkanswer 提取，
% 并与 2005/2005_数学真题_福建理科_答案要点.md 逐题比对后生成，勿手改（改源 MD 后重跑）。

\begin{gknotice}
\textbf{使用说明}：本节只给最终答案便于快速核对，完整推导见解析版。答案为编写者独立推导，\textbf{非官方标准答案}。
\end{gknotice}

\gkpart{选择题}
\begin{center}\begin{tabular}{ccccccccccccc}
\toprule
题号 & 1 & 2 & 3 & 4 & 5 & 6 & 7 & 8 & 9 & 10 & 11 & 12 \\
\midrule
答案 & B & A & A & C & D & C & A & D & B & D & C & C \\
\bottomrule
\end{tabular}\end{center}

\gkpart{填空题}
\begin{itemize}[leftmargin=2.6em]
  \item[\textbf{13.}] 240
  \item[\textbf{14.}] 0
  \item[\textbf{15.}] $\dfrac{1}{b-1}$
  \item[\textbf{16.}] 直线 $y = x$；$2^{x-3}$（或写关于 $x$ 轴对称，$g(x) = -3 - \log_2 x$ 等）
\end{itemize}

\gkpart{解答题}
\begin{itemize}[leftmargin=2.6em]
  \item[\textbf{17.}] (1) $\sin x - \cos x = -\dfrac{7}{5}$；(2) $\dfrac{3\sin^2\dfrac{x}{2} - 2\sin\dfrac{x}{2}\cos\dfrac{x}{2} + \cos^2\dfrac{x}{2}}{\tan x + \cot x} = -\dfrac{108}{125}$
  \item[\textbf{18.}] (1) $\mathbf{E}\xi = \dfrac{9}{10}$；(2) 这四次投球中至少一次命中的概率为 $\dfrac{91}{100}$
  \item[\textbf{19.}] (1) 函数解析式为 $f(x) = \dfrac{2x - 6}{x^2 + 3}$；(2) 单调递增区间为 $(3 - 2\sqrt{3}, 3 + 2\sqrt{3})$；单调递减区间为 $(-\infty, 3 - 2\sqrt{3})$ 和 $(3 + 2\sqrt{3}, +\infty)$
  \item[\textbf{20.}] (1) 证明见解析（由 $CB \perp AE$ 且 $BF \perp AE$ 导出 $AE \perp \text{\text{平面} } BCE$）；(2) 二面角 $B - AC - E$ 的大小为 $\arccos\dfrac{\sqrt{3}}{3}$；(3) 点 $D$ 到平面 $ACE$ 的距离为 $\dfrac{2\sqrt{3}}{3}$
  \item[\textbf{21.}] (1) 椭圆 $C$ 的方程为 $\dfrac{x^2}{6} + \dfrac{y^2}{2} = 1$；(2) 存在，所求直线 $m$ 的方程为 $x = -2$
  \item[\textbf{22.}] (1) $a = -\dfrac{2}{3}$；(2) 证明见解析（由逆推映射性质知经有穷步必到达 0 从而数列终止）；(3) $a$ 的取值范围为 $(0, +\infty)$（即 $a > 0$）
\end{itemize}
